\documentclass[12pt,a4paper]{article}

% ── Packages ──────────────────────────────────────────────────
\usepackage[utf8]{inputenc}
\usepackage[T1]{fontenc}
\usepackage[english]{babel}
\usepackage{geometry}
\geometry{margin=2.5cm}
\usepackage{graphicx}
\usepackage{booktabs}
\usepackage{hyperref}
\usepackage{xcolor}
\usepackage{listings}
\usepackage{amsmath}
\usepackage{enumitem}
\usepackage{float}
\usepackage{caption}
\usepackage{fancyhdr}
\usepackage{titlesec}
\usepackage{tocloft}

% ── Colors & Styles ───────────────────────────────────────────
\definecolor{codeblue}{HTML}{1a73e8}
\definecolor{codegray}{HTML}{6c757d}
\definecolor{codebg}{HTML}{f8f9fa}
\definecolor{darkgreen}{HTML}{28a745}

\hypersetup{
    colorlinks=true,
    linkcolor=codeblue,
    urlcolor=codeblue,
    citecolor=codeblue,
}

\lstset{
    backgroundcolor=\color{codebg},
    basicstyle=\ttfamily\small,
    breaklines=true,
    frame=single,
    rulecolor=\color{codegray},
    keywordstyle=\color{codeblue}\bfseries,
    commentstyle=\color{codegray},
    stringstyle=\color{darkgreen},
    showstringspaces=false,
    tabsize=4,
}

% ── Header/Footer ─────────────────────────────────────────────
\pagestyle{fancy}
\fancyhf{}
\fancyhead[L]{\small Binary Classification of Algae}
\fancyhead[R]{\small University Project}
\fancyfoot[C]{\thepage}

% ══════════════════════════════════════════════════════════════
%                        DOCUMENT
% ══════════════════════════════════════════════════════════════
\begin{document}

% ── Title Page ────────────────────────────────────────────────
\begin{titlepage}
    \centering
    \vspace*{2cm}
    {\Huge\bfseries Binary Image Classification\\of Marine Algae\\[0.5cm]}
    {\LARGE Using Transfer Learning with MobileNetV2\\[1.5cm]}
    {\Large\itshape Classifying \textbf{Alma} and \textbf{Glaci} Algae Species\\from Beach Photographs in Agadir, Morocco\\[2cm]}
    {\large University Project\\[0.5cm]}
    {\large\today\\[3cm]}
    \vfill
    {\normalsize Built with TensorFlow, Keras \& Streamlit}
\end{titlepage}

% ── Table of Contents ─────────────────────────────────────────
\newpage
\tableofcontents
\newpage

% ══════════════════════════════════════════════════════════════
\section{Introduction}
% ══════════════════════════════════════════════════════════════

This project implements an end-to-end \textbf{binary image classification} system that distinguishes between two types of marine algae---\textbf{Alma} and \textbf{Glaci}---found on the beaches of Agadir, Morocco. The system leverages \textbf{transfer learning} with the MobileNetV2 architecture pre-trained on ImageNet, combined with a custom classification head trained on our algae dataset. A Streamlit web application provides a user-friendly interface for real-time inference.

\subsection{Objectives}

\begin{enumerate}
    \item Collect and curate a photographic dataset of two algae types from Agadir beaches.
    \item Build an accurate deep learning classifier using transfer learning.
    \item Deploy the model through an interactive web interface.
    \item Document the biological context and technical methodology.
\end{enumerate}

\subsection{Key Technologies}

\begin{itemize}
    \item \textbf{TensorFlow / Keras} --- Deep learning framework
    \item \textbf{MobileNetV2} --- Pre-trained backbone (ImageNet weights)
    \item \textbf{Streamlit} --- Web-based inference UI
    \item \textbf{scikit-learn} --- Evaluation metrics
    \item \textbf{Matplotlib} --- Learning curve visualization
\end{itemize}

% ══════════════════════════════════════════════════════════════
\section{Biological Context}
% ══════════════════════════════════════════════════════════════

\subsection{Marine Algae Overview}

Marine algae (seaweeds) are photosynthetic organisms found in coastal and oceanic environments worldwide. They play a critical role in marine ecosystems by:

\begin{itemize}
    \item \textbf{Producing oxygen} --- Marine algae contribute an estimated 50\% of the world's oxygen through photosynthesis.
    \item \textbf{Carbon fixation} --- They absorb CO\textsubscript{2} from the atmosphere, acting as a carbon sink.
    \item \textbf{Habitat creation} --- Algae beds provide shelter, breeding grounds, and food for marine fauna.
    \item \textbf{Nutrient cycling} --- They recycle nutrients within marine food webs.
    \item \textbf{Economic value} --- Algae are used in food, cosmetics, pharmaceuticals, fertilizers, and biofuels.
\end{itemize}

Marine algae are broadly classified into three major groups based on their pigmentation:

\begin{table}[H]
    \centering
    \caption{Major groups of marine algae}
    \begin{tabular}{@{}llll@{}}
        \toprule
        \textbf{Group} & \textbf{Common Name} & \textbf{Primary Pigments} & \textbf{Example Genera} \\
        \midrule
        Chlorophyta  & Green algae  & Chlorophyll a, b      & \textit{Ulva}, \textit{Codium} \\
        Phaeophyta   & Brown algae  & Fucoxanthin           & \textit{Fucus}, \textit{Sargassum} \\
        Rhodophyta   & Red algae    & Phycoerythrin         & \textit{Gracilaria}, \textit{Gelidium} \\
        \bottomrule
    \end{tabular}
\end{table}

\subsection{Algae in Agadir, Morocco}

Agadir is a coastal city in southwestern Morocco, situated along the Atlantic Ocean. The Agadir coastline is characterized by:

\begin{itemize}
    \item A long sandy beach stretching over 10 km
    \item Temperate to warm Atlantic waters influenced by the Canary Current
    \item Rocky intertidal zones rich in biodiversity
    \item Seasonal upwelling that brings nutrient-rich deep water to the surface, promoting algal growth
\end{itemize}

The Moroccan Atlantic coast is known for its diverse marine flora. The region of Souss-Massa (where Agadir is located) hosts numerous species of macroalgae. Algae are commonly found attached to rocks in the intertidal zone or washed ashore on beaches.

\subsection{The Two Study Species}

This project focuses on two types of algae collected from the beaches of Agadir:

\subsubsection{Alma}

\textbf{Alma} refers to a type of algae commonly encountered on the Agadir shoreline. In the local context, ``alma'' is a vernacular name used to describe certain species of edible or commonly recognized seaweed. Key characteristics:

\begin{itemize}
    \item Found in the intertidal and subtidal zones of rocky shores
    \item Typically exhibits distinctive coloration and morphological features
    \item Can be distinguished visually from other algae types by its shape, texture, and color
    \item 300 photographs were collected for this class
\end{itemize}

\subsubsection{Glaci}

\textbf{Glaci} represents the second type of algae studied in this project. Key characteristics:

\begin{itemize}
    \item Also found along the Agadir coastline
    \item Exhibits different visual characteristics compared to Alma
    \item Distinct morphology (shape, branching pattern, texture, and/or coloration)
    \item 375 photographs were collected for this class
\end{itemize}

\subsection{Why Classify Algae Automatically?}

Accurate identification of algae species is important for:

\begin{enumerate}
    \item \textbf{Ecological monitoring} --- Tracking species distribution and abundance over time helps detect environmental changes.
    \item \textbf{Biodiversity assessment} --- Cataloguing algal diversity in Moroccan coastal waters.
    \item \textbf{Water quality indication} --- Certain algae thrive in polluted or nutrient-enriched waters; their presence can indicate water quality.
    \item \textbf{Commercial applications} --- Identifying commercially valuable species used for agar extraction, food, cosmetics, or biofuels.
    \item \textbf{Research acceleration} --- Automating identification frees researchers from tedious manual sorting.
\end{enumerate}

Traditional algae identification requires trained marine biologists who examine morphological features under magnification. This process is:
\begin{itemize}
    \item Time-consuming
    \item Dependent on expert availability
    \item Subject to inter-observer variability
\end{itemize}

An automated deep learning classifier addresses all of these limitations.

% ══════════════════════════════════════════════════════════════
\section{Dataset}
% ══════════════════════════════════════════════════════════════

\subsection{Data Collection}

Photographs were captured directly on the beaches of Agadir, Morocco using a standard camera. The images are natural photographs taken in real-world conditions (variable lighting, backgrounds, angles).

\subsection{Raw Data}

\begin{table}[H]
    \centering
    \caption{Raw dataset composition}
    \begin{tabular}{@{}lrl@{}}
        \toprule
        \textbf{Class} & \textbf{Images} & \textbf{Format} \\
        \midrule
        Alma    & 300 & JPG \\
        Glaci   & 375 & JPG \\
        \midrule
        \textbf{Total} & \textbf{675} & \\
        \bottomrule
    \end{tabular}
\end{table}

The raw data contains duplicate images (files with ``Copie'' in the filename). After deduplication, approximately \textbf{269 unique images} remain.

\subsection{Data Split}

The unique images were divided into three non-overlapping subsets:

\begin{table}[H]
    \centering
    \caption{Train / Validation / Test split}
    \begin{tabular}{@{}lrrrl@{}}
        \toprule
        \textbf{Split} & \textbf{Alma} & \textbf{Glaci} & \textbf{Total} & \textbf{Purpose} \\
        \midrule
        Train      & 105 & 92 & 197 & Weight optimization \\
        Validation & 19  & 17 & 36  & Hyperparameter monitoring \\
        Test       & 19  & 17 & 36  & Final unbiased evaluation \\
        \bottomrule
    \end{tabular}
\end{table}

\subsection{Preprocessing Pipeline}

At load time, each image undergoes:

\begin{enumerate}
    \item \textbf{Resizing} --- All images resized to $96 \times 96$ pixels.
    \item \textbf{Batching} --- Grouped into mini-batches of 32 images.
    \item \textbf{Normalization} --- Pixel values scaled from $[0, 255]$ to $[-1, 1]$ via MobileNetV2 preprocessing.
    \item \textbf{Prefetching} --- \texttt{tf.data.AUTOTUNE} overlaps loading with training.
    \item \textbf{Error handling} --- Corrupted files are skipped via \texttt{.ignore\_errors()}.
\end{enumerate}

% ══════════════════════════════════════════════════════════════
\section{Model Architecture}
% ══════════════════════════════════════════════════════════════

\subsection{Transfer Learning Approach}

Training a convolutional neural network from scratch requires large datasets (typically tens of thousands of images). With only $\sim$270 unique images, we employ \textbf{transfer learning}: reusing a model pre-trained on ImageNet (1.2 million images, 1000 classes) and adapting it to our binary algae classification task.

\subsection{Architecture Overview}

The model consists of two parts:

\begin{enumerate}
    \item \textbf{Feature Extractor (Frozen)} --- MobileNetV2 backbone with ImageNet weights.
    \item \textbf{Classification Head (Trainable)} --- Custom layers for binary output.
\end{enumerate}

\begin{table}[H]
    \centering
    \caption{Model architecture layer by layer}
    \begin{tabular}{@{}llll@{}}
        \toprule
        \textbf{Layer} & \textbf{Output Shape} & \textbf{Parameters} & \textbf{Trainable} \\
        \midrule
        Input                   & $(96, 96, 3)$         & 0         & --- \\
        MobileNetV2 Preprocess  & $(96, 96, 3)$         & 0         & --- \\
        MobileNetV2 Backbone    & $(3, 3, 1280)$        & $\sim$2.2M & \textbf{No} (frozen) \\
        Global Average Pooling  & $(1280)$              & 0         & --- \\
        Dropout (0.2)           & $(1280)$              & 0         & --- \\
        Dense (sigmoid)         & $(1)$                 & 1,281     & \textbf{Yes} \\
        \midrule
        \multicolumn{2}{@{}l}{\textbf{Total trainable parameters}} & \textbf{1,281} & \\
        \bottomrule
    \end{tabular}
\end{table}

\subsection{Why MobileNetV2?}

MobileNetV2 was chosen for the following reasons:

\begin{itemize}
    \item \textbf{Lightweight} --- Only $\sim$3.4M parameters total, enabling fast inference.
    \item \textbf{Pre-trained on ImageNet} --- Rich visual feature representations.
    \item \textbf{Depthwise separable convolutions} --- Efficient computation.
    \item \textbf{Inverted residual blocks} --- Better gradient flow during training.
    \item \textbf{Flexible input sizes} --- Works well at $96 \times 96$.
\end{itemize}

\begin{table}[H]
    \centering
    \caption{Comparison with other backbone architectures}
    \begin{tabular}{@{}lrl@{}}
        \toprule
        \textbf{Model} & \textbf{Parameters} & \textbf{Inference Speed} \\
        \midrule
        \textbf{MobileNetV2} & \textbf{3.4M} & \textbf{Fast} \\
        ResNet-50     & 25.6M  & Medium \\
        VGG-16        & 138M   & Slow \\
        \bottomrule
    \end{tabular}
\end{table}

\subsection{Layer Details}

\subsubsection{MobileNetV2 Preprocessing}

MobileNetV2 expects input pixels in the range $[-1, 1]$. The built-in preprocessing function applies:

$$x_{\text{normalized}} = \frac{x}{127.5} - 1$$

\subsubsection{MobileNetV2 Backbone (Frozen)}

\begin{itemize}
    \item Pre-trained on ImageNet (1.2M images, 1000 classes).
    \item \texttt{include\_top=False} removes the original 1000-class head.
    \item \texttt{base.trainable = False} freezes all backbone weights.
    \item Acts as a \textbf{fixed feature extractor}.
    \item Outputs a 3D feature map of shape $(3, 3, 1280)$ for $96 \times 96$ input.
\end{itemize}

\subsubsection{Global Average Pooling 2D}

Reduces the spatial dimensions by averaging over the $3 \times 3$ spatial grid:

$$\text{GAP}(F)_c = \frac{1}{H \times W} \sum_{i=1}^{H} \sum_{j=1}^{W} F_{i,j,c}$$

This produces a 1280-dimensional feature vector.

\subsubsection{Dropout (rate = 0.2)}

During training, randomly sets 20\% of activations to zero. This regularization technique helps prevent overfitting, which is especially important with our small dataset.

\subsubsection{Dense Layer (1 unit, Sigmoid)}

The final layer outputs a single value passed through the sigmoid function:

$$\sigma(z) = \frac{1}{1 + e^{-z}}$$

The output $p \in [0, 1]$ represents the probability of the input belonging to class ``glaci'':
\begin{itemize}
    \item $p \geq 0.5 \Rightarrow$ \textbf{Glaci}
    \item $p < 0.5 \Rightarrow$ \textbf{Alma}
\end{itemize}

% ══════════════════════════════════════════════════════════════
\section{Training Pipeline}
% ══════════════════════════════════════════════════════════════

\subsection{Training Configuration}

\begin{table}[H]
    \centering
    \caption{Configurable training parameters}
    \begin{tabular}{@{}lll@{}}
        \toprule
        \textbf{Parameter} & \textbf{Default} & \textbf{CLI Flag} \\
        \midrule
        Data directory  & \texttt{data\_split}                    & \texttt{--data-dir} \\
        Image size      & $96 \times 96$                          & \texttt{--img-size} \\
        Batch size      & 32                                      & \texttt{--batch-size} \\
        Epochs          & 10                                      & \texttt{--epochs} \\
        Model output    & \texttt{artifacts/2-epoches-test.keras} & \texttt{--model-out} \\
        \bottomrule
    \end{tabular}
\end{table}

\subsection{Compilation Settings}

\begin{table}[H]
    \centering
    \caption{Model compilation configuration}
    \begin{tabular}{@{}lll@{}}
        \toprule
        \textbf{Setting} & \textbf{Value} & \textbf{Rationale} \\
        \midrule
        Optimizer & Adam                & Adaptive learning rate \\
        Loss      & Binary Crossentropy & Standard for binary classification \\
        Metric    & Accuracy            & Intuitive training monitor \\
        \bottomrule
    \end{tabular}
\end{table}

\subsection{Loss Function}

Binary crossentropy is defined as:

$$\mathcal{L} = -\frac{1}{N}\sum_{i=1}^{N}\left[y_i \log(\hat{y}_i) + (1-y_i)\log(1-\hat{y}_i)\right]$$

where $y_i \in \{0, 1\}$ is the true label and $\hat{y}_i$ is the predicted probability.

\subsection{Training Workflow}

The complete training pipeline (\texttt{main.py}) executes the following steps:

\begin{enumerate}
    \item Parse command-line arguments
    \item Load train / validation / test datasets from \texttt{data\_split/}
    \item Build MobileNetV2 model with custom classification head
    \item Compile with Adam optimizer and binary crossentropy loss
    \item Train the model, monitoring validation performance each epoch
    \item Save the trained model to \texttt{artifacts/*.keras}
    \item Plot and save learning curves (loss and accuracy)
    \item Evaluate on train / validation / test sets
    \item Save metrics to \texttt{artifacts/metrics.json}
    \item Save metadata (class names, image size) to \texttt{artifacts/meta.json}
    \item Automatically update \texttt{README.md} with latest test metrics
\end{enumerate}

\subsection{Artifacts Produced}

\begin{table}[H]
    \centering
    \caption{Files generated after training}
    \begin{tabular}{@{}ll@{}}
        \toprule
        \textbf{File} & \textbf{Description} \\
        \midrule
        \texttt{2-epoches-test.keras} & Serialized trained Keras model \\
        \texttt{meta.json}            & Class names and image size \\
        \texttt{metrics.json}         & All evaluation metrics \\
        \texttt{learning\_curves.png} & Loss \& accuracy plots \\
        \bottomrule
    \end{tabular}
\end{table}

% ══════════════════════════════════════════════════════════════
\section{Evaluation \& Results}
% ══════════════════════════════════════════════════════════════

\subsection{Metrics}

Four standard classification metrics are computed:

\begin{align}
    \text{Accuracy}  &= \frac{TP + TN}{TP + TN + FP + FN} \\[6pt]
    \text{Precision} &= \frac{TP}{TP + FP} \\[6pt]
    \text{Recall}    &= \frac{TP}{TP + FN} \\[6pt]
    \text{F1-Score}  &= 2 \cdot \frac{\text{Precision} \times \text{Recall}}{\text{Precision} + \text{Recall}}
\end{align}

\subsection{Results}

\begin{table}[H]
    \centering
    \caption{Model performance across all data splits}
    \begin{tabular}{@{}lrrrr@{}}
        \toprule
        \textbf{Split} & \textbf{Accuracy} & \textbf{Precision} & \textbf{Recall} & \textbf{F1-Score} \\
        \midrule
        Train      & 1.0000 & 1.0000 & 1.0000 & 1.0000 \\
        Validation & 1.0000 & 1.0000 & 1.0000 & 1.0000 \\
        Test       & 1.0000 & 1.0000 & 1.0000 & 1.0000 \\
        \bottomrule
    \end{tabular}
\end{table}

\subsection{Analysis}

The model achieves \textbf{perfect classification} (100\%) across all splits. This indicates:

\begin{itemize}
    \item The two algae types are \textbf{visually very distinct}---the morphological and color differences between Alma and Glaci are clear enough for the model to learn a perfect decision boundary.
    \item Transfer learning features from ImageNet generalize well to marine biological imagery.
    \item There are \textbf{no signs of overfitting} (train performance = validation performance = test performance).
    \item The frozen MobileNetV2 features are more than sufficient for this task.
\end{itemize}

% ══════════════════════════════════════════════════════════════
\section{Inference Application}
% ══════════════════════════════════════════════════════════════

\subsection{Streamlit Web App}

The project includes an interactive web application built with Streamlit (\texttt{app.py}) for real-time inference.

\subsection{Features}

\begin{itemize}
    \item Upload any JPG/PNG image
    \item Instant classification result with confidence percentage
    \item Raw probability distribution for both classes
    \item Display of latest training metrics
    \item Model caching for fast repeated predictions
\end{itemize}

\subsection{Inference Pipeline}

\begin{enumerate}
    \item User uploads an image through the web interface.
    \item Image is converted to RGB (if necessary).
    \item Image is resized to $96 \times 96$ pixels.
    \item Image is converted to a NumPy array with batch dimension: shape $(1, 96, 96, 3)$.
    \item Model performs a forward pass (MobileNetV2 preprocessing applied internally).
    \item Sigmoid output $p$ is thresholded at 0.5 to determine the predicted class.
    \item The class label and confidence ($\max(p, 1-p)$) are displayed.
\end{enumerate}

\subsection{Running the App}

\begin{lstlisting}[language=bash]
streamlit run app.py
\end{lstlisting}

Then open \texttt{http://localhost:8501} in a browser.

% ══════════════════════════════════════════════════════════════
\section{Project Structure}
% ══════════════════════════════════════════════════════════════

\begin{lstlisting}[basicstyle=\ttfamily\footnotesize]
binary-classification/
|-- main.py                   # Training script
|-- app.py                    # Streamlit inference app
|-- requirements.txt          # Python dependencies
|-- README.md                 # Project readme
|-- data/                     # Raw images
|   |-- alma/                 # 300 JPGs
|   |-- glaci/                # 375 JPGs
|-- data_split/               # Train/Val/Test
|   |-- train/
|   |   |-- alma/   (105)
|   |   |-- glaci/  (92)
|   |-- val/
|   |   |-- alma/   (19)
|   |   |-- glaci/  (17)
|   |-- test/
|       |-- alma/   (19)
|       |-- glaci/  (17)
|-- artifacts/                # Training outputs
    |-- 2-epoches-test.keras
    |-- meta.json
    |-- metrics.json
    |-- learning_curves.png
\end{lstlisting}

% ══════════════════════════════════════════════════════════════
\section{Reproduction Guide}
% ══════════════════════════════════════════════════════════════

\subsection{Prerequisites}

\begin{itemize}
    \item Python 3.8+
    \item pip package manager
\end{itemize}

\subsection{Step 1: Install Dependencies}

\begin{lstlisting}[language=bash]
pip install -r requirements.txt
\end{lstlisting}

Dependencies: \texttt{tensorflow-cpu}, \texttt{pillow}, \texttt{streamlit}, \texttt{scikit-learn}.

\subsection{Step 2: Prepare Data}

Place algae photographs into the class folders:

\begin{lstlisting}[language=bash]
data/alma/    # Alma algae images
data/glaci/   # Glaci algae images
\end{lstlisting}

Then split into \texttt{data\_split/train}, \texttt{data\_split/val}, and \texttt{data\_split/test}.

\subsection{Step 3: Train}

\begin{lstlisting}[language=bash]
python main.py
# Or with custom parameters:
python main.py --epochs 20 --img-size 128 --batch-size 16
\end{lstlisting}

\subsection{Step 4: Run Inference}

\begin{lstlisting}[language=bash]
streamlit run app.py
\end{lstlisting}

% ══════════════════════════════════════════════════════════════
\section{Design Decisions}
% ══════════════════════════════════════════════════════════════

\begin{table}[H]
    \centering
    \caption{Key technical decisions and their rationale}
    \begin{tabular}{@{}p{4cm}p{8cm}@{}}
        \toprule
        \textbf{Decision} & \textbf{Rationale} \\
        \midrule
        Transfer learning over training from scratch & Only $\sim$270 unique images; training from scratch would severely overfit. ImageNet features generalize well as generic visual representations. \\
        \addlinespace
        Frozen backbone & Small dataset means fine-tuning all layers risks overfitting. Freezing the backbone limits trainable parameters to $\sim$1,281. \\
        \addlinespace
        $96 \times 96$ input size & Faster training/inference with lower memory. Sufficient for visually distinct classes. \\
        \addlinespace
        Sigmoid + Binary Crossentropy & Single output neuron is more efficient than 2-class softmax for binary tasks. \\
        \addlinespace
        Adam optimizer & Adaptive learning rate; works well out of the box without tuning. \\
        \addlinespace
        MobileNetV2 & Best trade-off between parameter count and feature quality for lightweight deployment. \\
        \bottomrule
    \end{tabular}
\end{table}

% ══════════════════════════════════════════════════════════════
\section{Potential Improvements}
% ══════════════════════════════════════════════════════════════

\begin{table}[H]
    \centering
    \caption{Possible future improvements}
    \begin{tabular}{@{}p{4.5cm}p{6.5cm}l@{}}
        \toprule
        \textbf{Improvement} & \textbf{Expected Benefit} & \textbf{Effort} \\
        \midrule
        Data Augmentation        & More robust to lighting/angle variations & Low \\
        Fine-Tuning Top Layers   & Better domain-specific features          & Low \\
        Larger Image Size        & Capture finer morphological details      & Low \\
        Early Stopping           & Prevent overfitting, save time           & Low \\
        Learning Rate Scheduling & Better convergence                       & Medium \\
        Grad-CAM Visualization   & Explainability---see what the model focuses on & Medium \\
        Multi-class Extension    & Classify more algae species              & Medium \\
        \bottomrule
    \end{tabular}
\end{table}

% ══════════════════════════════════════════════════════════════
\section{Conclusion}
% ══════════════════════════════════════════════════════════════

This project demonstrates that \textbf{transfer learning with MobileNetV2} is a highly effective approach for binary classification of marine algae, even with a very small dataset ($\sim$270 images). By reusing features learned from ImageNet and training only a lightweight classification head (1,281 parameters), we achieved \textbf{100\% accuracy} across all evaluation splits.

The combination of:
\begin{itemize}
    \item A pre-trained backbone for feature extraction
    \item A simple classification head
    \item A Streamlit web interface for deployment
\end{itemize}

\noindent makes this system practical, efficient, and accessible. The approach can be readily extended to classify additional algae species or applied to other biological image classification tasks.

\vspace{1cm}
\noindent\rule{\textwidth}{0.4pt}
\begin{center}
    \textit{End of Documentation}
\end{center}

\end{document}
